\chapter{Postpartum Cardiomyopathy}
\index{Postpartum Cardiomyopathy}
\label{sec:Postpartum Cardiomyopathy}

\section{Overview}
Peripartum cardiomyopathy (PPCM) is an idiopathic dilated cardiomyopathy presenting with left ventricular systolic dysfunction (LVEF <45\%) in the last month of pregnancy through five months postpartum, occurring in 1 in 1000--4000 pregnancies. The Disease mechanism involves oxidative stress, a 16-kDa prolactin cleavage fragment (vasoinhibin) that impairs cardiac microvascular function, and inflammation. This cardiovascular condition is a leading cause of morbidity and mortality globally. Risk increases with age, and the condition often requires lifelong management. This pregnancy-related condition requires careful monitoring to ensure the safety of both mother and baby. Early detection and management are essential. Cardiovascular disease encompasses conditions affecting the heart and blood vessels, representing a leading cause of morbidity and mortality worldwide. These conditions range from acute events like heart attack and stroke to chronic conditions requiring lifelong management. Atherosclerosis, the buildup of plaque in arterial walls, underlies many cardiovascular conditions. Risk increases with age, and the global burden continues to rise with aging populations and lifestyle changes.
\section{Causes}
Atherosclerosis begins with endothelial injury and dysfunction, leading to lipid accumulation and inflammatory cell infiltration in arterial walls. Plaque formation narrows vessels and reduces blood flow; plaque rupture can trigger acute thrombosis and vessel occlusion. Hemodynamic factors such as hypertension increase mechanical stress on vessel walls. Metabolic abnormalities including diabetes and dyslipidemia accelerate the atherosclerotic process. Genetic factors influence lipid metabolism, blood pressure regulation, and vascular health.
\section{Risk Factors}

\begin{itemize}[noitemsep]
  \item Advanced age
  \item Cigarette smoking
  \item Hypertension (high blood pressure)
  \item Diabetes mellitus
  \item Elevated LDL cholesterol and low HDL cholesterol
  \item Family history of premature cardiovascular disease
  \item Obesity (especially abdominal obesity)
  \item Physical inactivity
  \item Unhealthy diet (high in saturated fat, salt, processed foods)
\end{itemize}
\section{Signs and Symptoms}

\subsection{Common symptoms}
\begin{itemize}[noitemsep]
  \item You may experience shortness of breath, orthopnea, paroxysmal nocturnal shortness of breath, peripheral swelling, cough, and signs of left heart failure. Chest pain, pressure, or discomfort (angina) Shortness of breath with exertion or at rest Palpitations or irregular heartbeat Fatigue and reduced exercise tolerance Swelling in the legs, ankles, or feet (edema) Dizziness or lightheadedness Pain radiating to arm, jaw, back, or neck Nausea and indigestion Cold sweats Sudden weakness or numbness (stroke symptoms)
  \item Chest pain, pressure, or discomfort (angina)
  \item Shortness of breath with exertion or at rest
  \item Palpitations or irregular heartbeat
  \item Fatigue and reduced exercise tolerance
  \item Swelling in the legs, ankles, or feet (edema)
  \item Dizziness or lightheadedness
  \item Pain radiating to arm, jaw, back, or neck
  \item Nausea and indigestion
  \item Cold sweats
\end{itemize}

\subsection{Emergency warning signs}
\begin{itemize}[noitemsep]
  \item Severe or worsening symptoms of postpartum cardiomyopathy require immediate medical evaluation
\end{itemize}
\section{Progression and Stages}
Cardiovascular disease develops gradually over years, often beginning with asymptomatic risk factors. Early stages involve endothelial dysfunction and fatty streak formation in arterial walls. Progressive plaque buildup leads to hemodynamically significant stenosis causing symptoms with exertion. Advanced disease involves unstable plaques prone to rupture, causing acute coronary syndromes. End-stage disease results in chronic heart failure or critical limb ischemia.
\section{Complications}
\begin{itemize}[noitemsep]
  \item Myocardial infarction (heart attack)
  \item Heart failure with reduced or preserved ejection fraction
  \item Stroke from embolic or thrombotic events
  \item Arrhythmias including atrial fibrillation and ventricular tachycardia
  \item Peripheral artery disease causing limb ischemia
  \item Aortic aneurysm and dissection
  \item Chronic kidney disease from reduced perfusion
\end{itemize}
\section{Diagnosis}
Doctors diagnose this using echocardiogram showing left ventricular dilation and reduced ejection fraction. Cardiovascular diagnosis relies on a combination of clinical history, physical examination, electrocardiography, imaging (echocardiography, CT angiography), and laboratory markers (troponin, BNP). History and physical examination including blood pressure measurement Electrocardiogram (ECG/EKG) to assess heart rhythm and ischemia Echocardiogram to evaluate cardiac structure and function Stress testing (exercise or pharmacologic) for ischemic evaluation Cardiac biomarkers (troponin, CK-MB, BNP) Lipid profile and glucose testing Coronary angiography or CT angiography for coronary anatomy Holter monitoring for arrhythmia detection Cardiac MRI for detailed structural assessment Ankle-brachial index for peripheral artery disease

\begin{itemize}[noitemsep]
  \item History and physical examination including blood pressure measurement
  \item Electrocardiogram (ECG/EKG) to assess heart rhythm and ischemia
  \item Echocardiogram to evaluate cardiac structure and function
  \item Stress testing (exercise or pharmacologic) for ischemic evaluation
  \item Cardiac biomarkers (troponin, CK-MB, BNP)
  \item Lipid profile and glucose testing
  \item Coronary angiography or CT angiography for coronary anatomy
  \item Holter monitoring for arrhythmia detection
\end{itemize}
\section{Prevention}
\begin{itemize}[noitemsep]
  \item Maintain a heart-healthy diet low in saturated fat and sodium
  \item Engage in regular aerobic exercise (150 minutes per week)
  \item Avoid tobacco products and limit alcohol
  \item Control blood pressure, cholesterol, and blood glucose
  \item Maintain a healthy body weight
  \item Manage stress through relaxation techniques
  \item Take prescribed preventive medications (statins, aspirin)
\end{itemize}
\section{Treatment Options}
\begin{itemize}[noitemsep]
  \item Treatment typically involves standard heart failure therapy adjusted for pregnancy and lactation
  \item Bromocriptine may be used in some centers.
  \item Lifestyle modifications: heart-healthy diet, regular exercise, smoking cessation
  \item Antihypertensive medications (ACE inhibitors, ARBs, beta-blockers, diuretics)
  \item Lipid-lowering therapy (statins, ezetimibe, PCSK9 inhibitors)
  \item Antiplatelet therapy (aspirin, clopidogrel)
  \item Anticoagulation for atrial fibrillation and thromboembolic risk
  \item Nitrates and antianginal medications
  \item Revascularization: percutaneous coronary intervention (stenting)
  \item Coronary artery bypass grafting for severe disease
  \item Cardiac rehabilitation programs
\end{itemize}
\section{Living with It}
\begin{itemize}[noitemsep]
  \item Take all medications exactly as prescribed
  \item Monitor your blood pressure and weight regularly
  \item Follow a heart-healthy diet and stay physically active
  \item Attend cardiac rehabilitation if recommended
  \item Recognize warning signs of heart attack and stroke
  \item Keep all follow-up appointments with your cardiologist
  \item Manage stress through relaxation and mindfulness
  \item Carry a list of your medications and dosages
\end{itemize}
\section{Outlook}
Recovery varies: 50--60\% recover LVEF to $\ge$50\% within 6 months. With proper management and lifestyle modifications, many patients maintain good quality of life. Long-term outcomes depend on the extent of disease, adherence to treatment, and control of risk factors. Outcomes depend on the extent of disease, adherence to treatment, and control of risk factors. Early intervention for acute events significantly improves survival. Regular monitoring and medication adherence are essential for preventing disease progression. Advances in interventional cardiology continue to improve outcomes for complex cases.
